\documentclass[aspectratio=169,10pt,english]{beamer}

\input{../../../_preambulo_beamer.tex}
\input{../../../_comandos_beamer.tex}

\selectlanguage{english}

\title{MA1001B - Statistical Modeling for Decision Making}
\subtitle{Lesson 46: ANOVA Assumptions and Diagnostics}
\author{}
\institute{Tecnol\'ogico de Monterrey}
\date{}

\begin{document}

\begin{frame}[plain]
  \vspace{1.2cm}
  {\Large\bfseries MA1001B - Statistical Modeling for Decision Making\par}
  \vspace{0.7cm}
  {\large Lesson 46: ANOVA Assumptions and Diagnostics\par}
  \vspace{0.7cm}
  {\color{TecRojo}\rule{\textwidth}{1.2pt}}
  \vspace{0.8cm}

  MA1001B Faculty\\[0.3cm]
  Term\\[0.3cm]
  Tecnol\'ogico de Monterrey
\end{frame}

\begin{frame}{The Assumption Question}
  \begin{alertblock}{Guiding question}
    After a significant $F$ and Tukey follow-up, what must still be true of
    residuals, spreads, and the sampling design?
  \end{alertblock}

  \vspace{0.6em}
  By the end of this lesson, you can:
  \begin{itemize}
    \item form residuals $e_{ij}=y_{ij}-\bar y_i$ and inspect their histogram;
    \item use Levene's test and residual-versus-fitted plots;
    \item treat independence as a design claim;
    \item explain why equal $n$ does not rescue strong heteroscedasticity.
  \end{itemize}

  \vfill
  \scriptsize All cycle times used today are synthetic. No real company data are used.
\end{frame}

\begin{frame}{Three Assumptions, One Residual}
  \small
  \begin{center}
    \begin{tabular}{ll}
      \toprule
      \textbf{Assumption} & \textbf{What we check} \\
      \midrule
      Normality & Histogram of $e_{ij}=y_{ij}-\bar y_i$ \\
      Equal variances & Levene and residual versus fitted \\
      Independence & Randomized stations, not a plot \\
      \bottomrule
    \end{tabular}
  \end{center}

  \vspace{0.5em}
  \begin{exampleblock}{Balanced packing experiment}
    $n=12$ stations in each method. Residual mean is $0.000000$ by
    construction inside each group.
  \end{exampleblock}
\end{frame}

\begin{frame}[fragile]{Step 1: Residual Normality}
  \begin{columns}[T]
    \begin{column}{0.40\textwidth}
      \small
      \begin{lstlisting}
e = y - group_mean
W, p = shapiro(e)
      \end{lstlisting}

      \begin{block}{Verified output}
        Residual SD $2.630914$\\
        Shapiro $W=0.977891$\\
        $p=0.673807$
      \end{block}
    \end{column}
    \begin{column}{0.60\textwidth}
      \centering
      \includegraphics[width=\textwidth]{../figures/anova_diagnostics_01_residual_normality.png}
      \vspace{0.2em}
      \scriptsize Residual histogram from Step 1
    \end{column}
  \end{columns}
\end{frame}

\begin{frame}{Equal Spread Is a Second, Separate Check}
  \small
  Sample SDs: $3.067311$ (Standard), $2.353358$ (Guided), $2.660249$
  (Automated).
  \[
    \text{Levene }W=0.457821,\qquad p=0.636619.
  \]
  Residual versus fitted should not fan out. Independence is inherited from
  the random assignment of stations: a residual plot cannot prove it.
\end{frame}

\begin{frame}[fragile]{Step 2: Levene and Residual versus Fitted}
  \begin{columns}[T]
    \begin{column}{0.40\textwidth}
      \small
      \begin{lstlisting}
W, p = levene(s, g, a)
e = y - fitted
      \end{lstlisting}

      \begin{block}{Verified output}
        Levene $p=0.636619$\\
        SDs $3.067$, $2.353$, $2.660$\\
        No evidence of unequal spread
      \end{block}
    \end{column}
    \begin{column}{0.60\textwidth}
      \centering
      \includegraphics[width=\textwidth]{../figures/anova_diagnostics_02_levene_homoscedasticity.png}
      \vspace{0.2em}
      \scriptsize Residual-versus-fitted plot from Step 2
    \end{column}
  \end{columns}
\end{frame}

\begin{frame}{Step 3: Equal $n$ Does Not Rescue Strong Heteroscedasticity}
  \begin{columns}[T]
    \begin{column}{0.42\textwidth}
      \small
      Still $n=12$ per method. SDs $1.534$, $1.324$, $9.976$.

      \vspace{0.4em}
      Levene $p=2.455155\times10^{-4}$. ANOVA $F=9.678598$,
      $p=4.924585\times10^{-4}$.

      \vspace{0.5em}
      \textbf{Limit:} equal-$n$ ANOVA is not immune to strong variance
      inequality. Lesson 47 moves from group means to a fitted line.
    \end{column}
    \begin{column}{0.58\textwidth}
      \centering
      \includegraphics[width=\textwidth]{../figures/anova_diagnostics_03_variance_inequality_limit.png}
      \vspace{0.2em}
      \scriptsize Heteroscedastic equal-$n$ example from Step 3
    \end{column}
  \end{columns}
\end{frame}

\begin{frame}{Concept Check}
  \begin{enumerate}
    \item Why is the residual mean exactly 0 inside each packing method?
    \item What does Levene $p=0.636619$ say, and what does it not say?
    \item Why can a residual plot not prove independence?
    \item In Step 3, $n=12$ in every group. Why is the significant $F$ still
      untrustworthy?
  \end{enumerate}

  \vspace{0.6em}
  \begin{block}{Connection to Lesson 47}
    The session can continue directly into simple linear regression, least
    squares, and $R^2$.
  \end{block}

  \vfill
  \scriptsize Source: OpenStax, \textit{Introductory Business Statistics 2e},
  Sections 12.2 and 12.4. CC BY-NC-SA.
\end{frame}

\appendix



\end{document}
